\subsection{Grothendieck classes and semiperfectness: Problem 21.60}
\label{cand:21.60}
\problemstatement{21.60}

Let \(G=C_3\rtimes C_4=\langle a,b:a^3=b^4=1,\
bab^{-1}=a^{-1}\rangle\), \(R=\Z_{(2)}\), and
\(k=\F_2\). Every simple \(kG\)-module is the reduction
of a simple rational module, but \(RG\) is not
semiperfect. Thus the nonnegative-combination criterion
in the question does not suffice.

The central involution \(b^2\) acts trivially on a
simple \(kG\)-module: its difference from the identity
is nilpotent with a nonzero invariant kernel.
Every such module factors through \(S_3\). In \(kS_3\),
the central idempotent \(e=1+a+a^2\) cuts out \(kC_2\),
whose only simple module is \(k\). The complementary
four-dimensional algebra is \(M_2(k)\), via
\[
 A=\begin{pmatrix}0&1\\1&1\end{pmatrix},\qquad
 B=\begin{pmatrix}1&1\\0&1\end{pmatrix}.
\]
These satisfy the relations, and \(I,A,B,AB\) are
linearly independent, proving the identification.
Its only simple module is the natural module \(V\).
This lifts to the rational representation with matrices
\[
 A_{\Q}=\begin{pmatrix}0&-1\\1&-1\end{pmatrix},\qquad
 B_{\Q}=\begin{pmatrix}1&-1\\0&-1\end{pmatrix}.
\]
The \(a\)-matrix has irreducible polynomial
\(X^2+X+1\), proving rational simplicity, and the
stable lattice \(R^2\) reduces to \(V\).
The trivial module lifts trivially. Hence every
finite-dimensional \(kG\)-module has Grothendieck
class a nonnegative combination of reductions of
simple rational modules, by a composition series.

For the obstruction to semiperfectness, let
\(E=\Q(\omega)\), \(\omega^2+\omega+1=0\), and
form the quaternion algebra
\[
 D=E+Ej,\qquad j^2=-1,\qquad ju=\bar u j.
\]
The norm of \(u+vj\) is \(u\bar u+v\bar v\).
For \(u=x+y\omega\), its summand is
\(x^2-xy+y^2\); the norm is positive on nonzero
elements. The conjugate divided by the norm gives
an inverse, so \(D\) is a division algebra.
The order \(O=R[\omega]+R[\omega]j\) is a quotient
of \(RG\) by \(a\mapsto\omega,b\mapsto j\).
Its reduction is \(M_2(k)\), since \(\omega,j\)
map to the matrices \(A,B\) above and their four
basis images are independent.

Moreover \(J(O)=2O\). For every \(y\in O\), the norm
of \(1-2y\) is congruent to one modulo \(2R\),
and quaternion conjugation preserves \(O\);
the inverse formula shows \(1-2y\) is a unit.
The radical unit criterion gives \(2O\subseteq J(O)\),
and the semisimplicity of \(O/2O\) gives the reverse
inclusion. The nontrivial idempotents of \(M_2(k)\)
cannot lift to \(O\), because a subring of a division
ring has only idempotents zero and one.
Thus \(O\) is not semiperfect.

A quotient of a semiperfect ring is semiperfect:
its quotient by its Jacobson radical is a quotient
of the original semisimple quotient; idempotents
lift first to that semisimple ring, then to the
original ring, and finally to the quotient.
Consequently \(RG\) is not semiperfect either.
The coefficient ring is the rational localization,
not the \(2\)-adic completion.

The question is Conjecture~9 in the retained version
of \cite{johnston-rumynin}. The example also conflicts
with its Proposition~3.3 as written. The extension-lifting
step there cannot hold for arbitrary chosen lattices:
with trivial \(C_2\)-actions,
\[
 \operatorname{Ext}^1_{RC_2}(R,R)=0,\qquad
 \operatorname{Ext}^1_{kC_2}(k,k)=k.
\]
Directly, an extension with fixed trivial submodule
and quotient has generator matrix
\(\left(\begin{smallmatrix}1&t\\0&1\end{smallmatrix}\right)\);
its square forces \(2t=0\) over \(R\), whereas
\(t=1\) gives a nonsplit extension over \(k\).
This concerns the cited version and step, and does
not import the proposition as a valid premise.
Source: \artifact{research/21.60-proof.md}.
